\chapter{Verified extensions and counterexamples}
\index{verified extension}\index{counterexample!cross-fan}\index{open problem!boundary ledger}
\label{chap:verified-extensions}

This chapter separates two provenance streams.  Its first-draft extensions
and counterexamples are restricted to the signatures that survived the
independent Rethlas audits recorded in the research-question ledger.  The
trace comparison and uniqueness statements, by contrast, are reconstructed
from \cite{DX24} and the complete trace chapter supplied with this project;
they are source imports, not new Rethlas-certified results.  Each theorem
below has a closed proof or source reconstruction, and each counterexample
identifies the exact failed implication.  Problems for which an audit found
only an idea or a missing approximation theorem are listed at the end and
are not inserted into the theorem narrative.

Verification is statement-specific. A receipt for a PL comparison on a
good strict nontrivially valued space does not authorize a theorem for a
trivially valued analytification, and a receipt for an abstract calibrated
fan does not construct a Berkovich analytic realization.  The historical
trace audit used an incomplete source packet and established only a weaker
\(\mathcal I\)-valued consequence.  It is superseded as a description of the
source frontier: the complete source proves that trace depends on the input
\(\mathcal I\)-class, has \(\mathcal I\)-good representatives, and is unique
as an output \(P\)-class.  The current \(P\)-valued formulation rests on that
source theorem, not on the obsolete receipt.  These firewalls determine the
wording of the chapter.

The seven outcomes are:
\begin{enumerate}
\item a sufficient CLD pluriharmonic-approximation criterion;
\item a cross-fan counterexample to unbalanced cellwise gluing;
\item the fixed-dual-complex case of the BJ Radon--Riesz question;
\item the exact \(P\)-class factorization and split trace solution;
\item a toric formula for Cao numerical dimension;
\item a no-go theorem for \(\mathcal I\)-invariant mass bodies; and
\item an equivalence among model, classical, and CLD-v2 positivity for
      real PL functions on boundaryless good strict spaces.
\end{enumerate}

\section{A sufficient CLD approximation criterion}
\label{sec:verified-cld-criterion}

Let \(U\) be an open subset of a good Berkovich analytic space on which
Chambert--Loir--Ducros forms and currents are defined. Let \(T\) be a
positive closed current on \(U\). Positive currents are of order zero,
so they extend continuously from smooth compactly supported test forms
to continuous ones. Consequently multiplication of \(T\) by a
continuous function is defined by
\[
                   \langle hT,\alpha\rangle
                       :=\langle T,h\alpha\rangle
\tag{69.1}\label{eq:verified-current-multiplication}
\]
for compactly supported smooth test forms \(\alpha\); the expression on
the right uses that order-zero extension, since \(h\alpha\) need not be
smooth.

\begin{theorem}\label{thm:verified-cld-criterion}
Let \(u,v\) be continuous psh functions on \(U\) satisfying
\[
                            \dprimeddprime u=\dprimeddprime v.
\tag{69.2}\label{eq:verified-equal-hessians}
\]
Put \(h=u-v\). Assume that there are smooth functions \(h_j\) on \(U\)
such that
\[
                h_j\longrightarrow h
          \quad\text{locally uniformly},
       \qquad
                \dprimeddprime h_j=0.
\tag{69.3}\label{eq:verified-pluriharmonic-approximation}
\]
Then
\[
                              \dprimeddprime(hT)=0.
\tag{69.4}\label{eq:verified-independence}
\]
\end{theorem}

\begin{proof}
Fix a compactly supported smooth test form \(\alpha\) of complementary
bidegree. Its support lies in a compact set \(K\Subset U\). Positivity
implies that \(T\) has finite mass on \(K\), hence
\[
 \bigl|\langle(h_j-h)T,\alpha\rangle\bigr|
 \leq
 \|h_j-h\|_{L^\infty(K)}
 \,\|T\|_K\,\|\alpha\|_{L^\infty(K)}.
\tag{69.5}\label{eq:verified-current-uniform}
\]
The right side tends to zero, so \(h_jT\to hT\) as currents.
Differentiation is continuous in the current topology. Since \(T\) is
closed and \(h_j\) is smooth,
\[
 \dprimeddprime(h_jT)
   =\dprimeddprime h_j\wedge T
   =0.
\tag{69.6}\label{eq:verified-smooth-product}
\]
Passing to the current limit gives
\[
 \dprimeddprime(hT)
   =\lim_j \dprimeddprime(h_jT)=0.
\]
\end{proof}

\begin{remark}[Sufficient, not automatic]
\label{rem:verified-cld-sufficient}
Equation \eqref{eq:verified-equal-hessians} does not produce the
approximants \((h_j)\). The theorem is therefore a conditional
independence statement. Existence of smooth pluriharmonic
approximants for every continuous distributionally pluriharmonic
function remains outside the result.
\end{remark}

\begin{corollary}\label{cor:verified-cld-smooth}
If \(h=u-v\) is itself smooth and \(\dprimeddprime h=0\), then
\(\dprimeddprime(hT)=0\).
\end{corollary}

\begin{proof}
Take the constant approximation \(h_j=h\) in
\Cref{thm:verified-cld-criterion}.
\end{proof}

\paragraph{One-dimensional check.}
Let \(U\) be a tropical interval, let \(T\) be its integration current, and
suppose \(u-v=ax+b\). Taking the constant sequence \(h_j=ax+b\) gives
\(\dprimeddprime((u-v)T)=0\). If the interval is replaced by a graph with a
vertex, the same calculation works on every open edge but says nothing
about the sum of outgoing slopes at the vertex. The next section isolates
exactly that missing gluing datum.

\section{Cell-interior rigidity and the gluing obstruction}
\label{sec:verified-cross-fan}

On the relative interior \(\sigma^\circ\) of a tropical cell, choose
affine coordinates \(x_1,\ldots,x_q\).

\begin{proposition}\label{prop:verified-cell-rigidity}
Let \(w\in L^1_{\mathrm{loc}}(\sigma^\circ)\). If
\[
                              \dprimeddprime[w]=0
\tag{69.7}\label{eq:verified-cell-harmonic}
\]
as a current on \(\sigma^\circ\), then \(w\) agrees almost everywhere
with an affine function. In particular, it has a unique continuous
affine representative.
\end{proposition}

\begin{proof}
In affine coordinates, the coefficients of \(\dprimeddprime[w]\) are the
distributional second derivatives \(\partial_i\partial_jw\). Thus
\[
                         \partial_i\partial_jw=0
\tag{69.8}\label{eq:verified-second-derivatives}
\]
for all \(i,j\). It follows that each first distributional derivative
\(\partial_iw\) is constant on the connected cell interior. Let those
constants be \(a_i\). Then the distributional gradient of
\[
                    w-\sum_{i=1}^qa_ix_i
\]
vanishes, so this difference is almost everywhere constant. The
resulting affine function is continuous. Two continuous functions that
agree almost everywhere on the cell agree everywhere, proving
uniqueness.
\end{proof}

The conclusion does not glue through a singular face without a
balancing law compatible with the current being multiplied.

\begin{proposition}\label{prop:verified-cross-fan}
Let
\[
 \Sigma
 =\mathbf R_{\geq0}e_1
   \cup\mathbf R_{\geq0}(-e_1)
   \cup\mathbf R_{\geq0}e_2
   \cup\mathbf R_{\geq0}(-e_2)
       \subset\mathbf R^2.
\tag{69.9}\label{eq:verified-cross-fan}
\]
Regard \(\Sigma\) as an abstract calibrated tropical fan and equip it first
with unit weights. Define on \(\Sigma\)
\[
                  u=2|x_1|,
       \qquad
                  v=|x_1|+|x_2|,
       \qquad
                  h=u-v.
\tag{69.10}\label{eq:verified-cross-functions}
\]
Then \(u\) and \(v\) are psh for the unit-weighted fan and
\(\dprimeddprime[u]=\dprimeddprime[v]=4\delta_0\). Now give each horizontal ray weight
\(2\) and each vertical ray weight \(1\), and let \(T\) be the
corresponding positive closed skeletal current. Despite the preceding
equality,
\[
                         \dprimeddprime(hT)=2\delta_0\neq0.
\tag{69.11}\label{eq:verified-cross-defect}
\]
\end{proposition}

\begin{proof}
First we verify psh positivity in the CLD sense, rather than infer it from
the outgoing Laplacian. Choose \(\varepsilon_j\downarrow0\) and define the
smooth ambient convex functions
\[
 u_j(x_1,x_2)=2\sqrt{x_1^2+\varepsilon_j},
 \qquad
 v_j(x_1,x_2)=\sqrt{x_1^2+\varepsilon_j}
              +\sqrt{x_2^2+\varepsilon_j}.
\]
Their restrictions to \(\Sigma\) are smooth psh functions and decrease
locally uniformly to \(u\) and \(v\), respectively. Closure of the CLD psh
class under decreasing limits therefore proves that \(u\) and \(v\) are
psh on the unit-weighted fan.

Both limit functions are convex on each ray. Their outgoing slopes are
\[
\begin{array}{c|cccc|c}
 &+e_1&-e_1&+e_2&-e_2&\text{sum}\\ \hline
u&2&2&0&0&4\\
v&1&1&1&1&4\\
h&1&1&-1&-1&0.
\end{array}
\tag{69.12}\label{eq:verified-cross-slopes}
\]
Thus the unit-weighted Laplacians of \(u\) and \(v\) are both
\(4\delta_0\), and the ordinary balancing defect of \(h\) at \(0\)
vanishes.
With the weights of \(T\), however, the outgoing defect of \(hT\) is
\[
             2\cdot1+2\cdot1+1\cdot(-1)+1\cdot(-1)=2.
\tag{69.13}\label{eq:verified-cross-weighted-sum}
\]
On the interior of every ray, \(h\) is affine, so there is no other
second-derivative contribution. Therefore
\[
                         \dprimeddprime(hT)=2\delta_0.
\]
The weighted fan defining \(T\) is balanced because opposite rays have equal
weights, and its weights are positive; hence \(T\) is positive and closed.
\end{proof}

\begin{remark}[Exact failed implication]
\label{rem:verified-cross-failed}
\Cref{prop:verified-cell-rigidity} remains true on each ray.
What fails is the passage from cellwise affinity plus an unweighted
balance to the balance required after multiplication by \(T\).
The obstruction is the change of face weights.
This is a counterexample on an abstract calibrated tropical fan. No
Berkovich analytic realization of this particular weighted configuration is
asserted here.
\end{remark}

An analytic realization would require more than drawing the same fan. One
would need a good analytic domain, a tropicalization chart carrying the four
rays, an analytic positive closed current whose tropical weights are
\((2,2,1,1)\), and compatibility of the current product with pushforward to
the fan. None of these data is supplied by
\eqref{eq:verified-cross-fan}. Accordingly the proposition disproves an
abstract tropical gluing principle, not an unqualified statement about all
Berkovich spaces.

\section{The fixed-dual-complex Radon--Riesz case}
\label{sec:verified-fixed-complex}

Work in the trivially valued projective setting of
\Cref{chap:bj-linear-growth}. Let \(X^{\lin}\) be the space of
linear-growth valuations. Its weak topology is induced from
\(X^{\An}\), and its strong topology is defined by the metric
\(d_{\infty,\omega}\) of \eqref{eq:bj-dinfty}. Write
\[
                  T_\omega(v)
                  =d_{\infty,\omega}(v,v_{\mathrm{triv}})
\tag{69.14}\label{eq:verified-exhaustion}
\]
for the standard exhaustion. Strong convergence implies weak
convergence and convergence of \(T_\omega\). Whether the converse holds
on all of \(X^{\lin}\) is Boucksom--Jonsson
Question~11.16 in \cite{BJ22}.

\begin{lemma}\label{lem:verified-compact-hausdorff}
Let \(K\subset X^{\lin}\) be compact for the strong topology.
Assume the weak topology is Hausdorff and is coarser than the strong
topology. Then the two induced topologies on \(K\) coincide.
\end{lemma}

\begin{proof}
The identity map
\[
     (K,\text{strong})\longrightarrow(K,\text{weak})
\tag{69.15}\label{eq:verified-identity-topologies}
\]
is a continuous bijection from a compact space to a Hausdorff space.
It is therefore a homeomorphism.
\end{proof}

\begin{theorem}\label{thm:verified-fixed-dual-complex}
Assume that the trivially valued ground field has characteristic zero and
that \(X\) is smooth and projective. Fix an SNC test configuration
\(\mathcal X\), and let
\(\Delta_{\mathcal X}\) be its dual complex, viewed as a compact
subset of \(X^{\lin}\). If a net
\((v_i)\subset\Delta_{\mathcal X}\) satisfies
\[
                        v_i\longrightarrow v
             \quad\text{weakly},
       \qquad
                        T_\omega(v_i)\longrightarrow T_\omega(v),
\tag{69.16}\label{eq:verified-radon-riesz-data}
\]
then \(v_i\to v\) strongly. In fact, weak convergence alone implies
strong convergence on this fixed dual complex.
\end{theorem}

\begin{proof}
The fixed dual complex is a finite compact polyhedral complex.
Its embedding into \((X^{\lin},d_{\infty,\omega})\) is Lipschitz by
the PL section formula \eqref{eq:bj-dinfty-sections} and the equi-Lipschitz
estimate of \cite[Theorem~A.4]{BJ22}. This is where characteristic zero and
smoothness enter. Its image is therefore strongly compact. The weak topology
is Hausdorff and coarser than the strong topology. Apply
\Cref{lem:verified-compact-hausdorff}. Thus weak and strong
topologies coincide on \(\Delta_{\mathcal X}\), so the first condition
in \eqref{eq:verified-radon-riesz-data} already gives the conclusion.
\end{proof}

\begin{corollary}\label{cor:verified-strongly-compact}
The same conclusion holds on every strongly compact subset of
\(X^{\lin}\).
\end{corollary}

\begin{proof}
Apply \Cref{lem:verified-compact-hausdorff} to that subset.
\end{proof}

\begin{remark}[No global conclusion]
\label{rem:verified-no-global-radon-riesz}
The compact-to-Hausdorff argument supplies no global compactness of
\(X^{\lin}\); in positive dimension, radial scaling already produces
strongly unbounded subsets. Comparability between
\(E^\vee(\delta_v)\) and \(T_\omega(v)\) does not identify their
continuity loci. Thus this argument does not answer the global question.
\end{remark}

\paragraph{Worked compact cell.}
If a fixed dual complex is an edge with affine coordinate \(s\in[0,1]\),
the PL section formula bounds
\(d_{\infty,\omega}(v_s,v_t)\) by \(C|s-t|\). Hence an ordinary convergent
net \(s_i\to s\) converges strongly. The theorem is the finite-complex
version of this calculation: finitely many Lipschitz charts replace the
single interval, and compactness prevents the Lipschitz constants from
escaping to infinity.

\section{Trace factorization and the split calculation}
\label{sec:verified-trace}

Return to a smooth compact K\"ahler manifold \(X\) and a closed
submanifold \(Y\). Recall the natural domain \(D_Y\) from
\eqref{eq:functor-trace-domain}.

\begin{theorem}\label{thm:verified-trace-factorization}
The natural trace assignment is
\[
 [\varphi]_{\mathcal I}\longmapsto
 [\Tr_Y(\varphi)]_P.
\tag{69.17}\label{eq:verified-trace-map}
\]
It is a well-defined map
\(D_Y/{\sim_{\mathcal I}}\to
\QPSH(\widetilde Y)/{\sim_P}\), and every output has an
\(\mathcal I\)-good representative.  Consequently its precomposition
with the quotient from \(P\)-classes factors as
\[
\begin{tikzcd}[column sep=large]
D_Y/{\sim_P}
 \arrow[r,two heads]
&
D_Y/{\sim_{\mathcal I}}
 \arrow[r]
&
\QPSH(\widetilde Y)/{\sim_P}.
\end{tikzcd}
\tag{69.18}\label{eq:verified-trace-factorization}
\]
Thus the source quotient is \(\mathcal I\)-theoretic, whereas the target
quotient is the finer \(P\)-quotient.  This asymmetry is the content of
the trace comparison and uniqueness theorems.
\end{theorem}

\begin{proof}
If \(\varphi\sim_{\mathcal I}\psi\), equality of multiplier-ideal
types gives equality of generic Lelong numbers at divisorial
valuations on every modification. If \(Y\) is a divisor, use
\(\ord_Y\); in higher codimension, use the exceptional divisor of the
blow-up of the smooth center \(Y\). In either case,
\[
                       \nu(\varphi,Y)=\nu(\psi,Y),
\tag{69.19}\label{eq:verified-trace-domain-saturation}
\]
so \(D_Y\) is \(\mathcal I\)-saturated. It is \(P\)-saturated because
\(P\)-equivalence implies \(\mathcal I\)-equivalence.

The trace comparison theorem gives
\[
 \varphi\preceq_{\mathcal I}\psi
       \Longrightarrow
 \Tr_Y(\varphi)\preceq_P\Tr_Y(\psi).
\tag{69.20}\label{eq:verified-trace-comparison}
\]
Applying this in both directions shows that
\(\mathcal I\)-equivalent inputs have \(P\)-equivalent traces.  The
trace uniqueness theorem also says that each representative is
\(\mathcal I\)-good.  Hence
the second arrow in \eqref{eq:verified-trace-factorization} is
well-defined and is precisely \eqref{eq:verified-trace-map}. Precomposing
with the quotient from \(P\)-classes proves the displayed factorization.
\end{proof}

\begin{theorem}\label{thm:verified-split-trace}
Let \(X=Y\times Z\) and
\[
                 \alpha=p_Y^*\alpha_Y+p_Z^*\alpha_Z,
\]
where both summands are K\"ahler classes. Let
\(S=\theta_Y+\ddc\psi\in\alpha_Y\) be a K\"ahler current and let
\(\eta\in\alpha_Z\) be a smooth K\"ahler form. Then
\[
                         T=p_Y^*S+p_Z^*\eta
\tag{69.21}\label{eq:verified-split-current}
\]
is K\"ahler and
\[
 \Tr_{Y\times\{z_0\}}(T)
       \sim_P P_{\alpha_Y}[\psi]_{\mathcal I}
       \sim_{\mathcal I}S.
\tag{69.22}\label{eq:verified-split-trace}
\]
If \(S\) is \(\mathcal I\)-good, its canonical representative
\(P_{\alpha_Y}[\psi]_{\mathcal I}\) is \(P\)-equivalent to \(S\), and
hence \(\Tr_{Y\times\{z_0\}}(T)\sim_P S\).
\end{theorem}

\begin{proof}
The pullback of \(S\) is strictly positive in horizontal directions,
and \(p_Z^*\eta\) is strictly positive in vertical directions, so
their sum is a K\"ahler current in \(\alpha\).

Choose a quasi-equisingular approximation \((\psi_j)\) of \(\psi\).
Since \(p_Y\) is smooth, local Fubini integration gives
\[
 \mathcal I(mp_Y^*\psi_j)
   =p_Y^{-1}\mathcal I(m\psi_j)\cdot\mathcal O_X
\tag{69.23}\label{eq:verified-smooth-pullback-ideals}
\]
for all \(m\), and likewise for \(\psi\). Thus
\((p_Y^*\psi_j)\) is quasi-equisingular for \(p_Y^*\psi\). Restriction
to \(Y\times\{z_0\}\) gives exactly \(\psi_j\), while the smooth
vertical term restricts to zero. The definition of trace therefore
gives
\[
 \Tr_{Y\times\{z_0\}}(T)
       \sim_P d_S\text{-}\lim_j\psi_j
       \sim_P P_{\alpha_Y}[\psi]_{\mathcal I}.
\tag{69.24}\label{eq:verified-split-limit}
\]
The last envelope is \(\mathcal I\)-equivalent to \(S\). If \(S\) is
\(\mathcal I\)-good, its \(P\)-class equals that of its canonical
\(\mathcal I\)-model envelope. This proves all assertions.
\end{proof}

\begin{remark}[Retraction form]
\label{rem:verified-retraction}
The same proof applies to a holomorphic submersion
\(r:X\to Y\) with section \(i\), provided
\(\alpha=r^*\beta+\gamma\) and \(\gamma\) has a smooth semipositive
representative \(\eta\) that is strictly positive on \(\ker(\rd r)\) and
satisfies \(i^*\eta=0\). Then \(T=r^*S+\eta\).
The transverse positivity and class splitting cannot be deleted.
\end{remark}

\section{Toric numerical dimension}
\label{sec:verified-toric-dimension}

Let \(X=X_\Sigma\) be a smooth projective toric variety of dimension
\(n\). Let \(D\) be an integral or rational big toric divisor and let
\(\theta\in c_1(D)\) be torus invariant. For a torus-invariant
\(\theta\)-psh potential \(\varphi\), put
\[
                         T=\theta+\ddc\varphi.
\]
Its Newton body is the compact convex set
\[
                         K=\Delta(T)
                          =\overline{\operatorname{dom}F_T^*}
                          \subset M_{\mathbf R}
\tag{69.25}\label{eq:verified-newton-body}
\]
defined from the convex weight on the dense torus.  The closure prevents
boundary slopes from being lost; compare
\Cref{def:ok-closed-newton-body}.

\begin{definition}[Cao numerical dimension]
\label{def:verified-cao-dimension}
After replacing \((D,T,\varphi,K)\) by
\((qD,qT,q\varphi,qK)\) for a common denominator \(q\), and retaining
the notation, regard \(\varphi\) as a singular semipositive metric on
an integral line bundle. Cao defines
\[
 \nu_{\mathrm{Cao}}(T)
   :=\max\bigl\{j\in\{0,\ldots,n\}:
       \langle T^j\rangle_{\mathrm{Cao}}\ne0\bigr\},
\]
where \(\langle T^j\rangle_{\mathrm{Cao}}\) is the cohomological
product constructed from quasi-equisingular approximations in
\cite[Definitions~3--4]{Cao14}. Multiplying \(T\) by a positive
integer does not change which of these classes vanish, so this also
defines the rational case.
\end{definition}

For later use, recall Cao's projective numerical criterion. If \(A\)
is very ample and \(\varepsilon>0\) is rational, put
\[
 S_A(\varepsilon)
 :=\varlimsup_{\substack{m\to\infty\\m\varepsilon\in\mathbf Z}}
 \frac{h^0\!\left(X,
   \mathcal O_X(mD+m\varepsilon A)\otimes
   \mathcal I(m\varphi)\right)}{m^n}.
\]
Then \cite[Proposition~4.4]{Cao14} says that
\[
 \lim_{\substack{\varepsilon\downarrow0\\
                  \varepsilon\in\mathbf Q}}
       \frac{\log S_A(\varepsilon)}{\log\varepsilon}
       =n-\nu_{\mathrm{Cao}}(T).
\tag{69.26}\label{eq:verified-cao-characterization}
\]

\begin{theorem}\label{thm:verified-toric-dimension}
Under the integral or rational hypotheses above,
\[
                          \nu_{\mathrm{Cao}}(T)=\dim K.
\tag{69.27}\label{eq:verified-toric-dimension}
\]
\end{theorem}

\begin{proof}
This is the integral/rational case of
\Cref{thm:ok-toric-real-products}.  We retain the following independent
section-counting proof because it explains how multiplier ideals detect the
same convex body when a graded section algebra exists.

Choose a very ample integral toric divisor \(A\) with full-dimensional
polytope \(Q\). The convex dictionary is additive under addition of
toric currents, hence
\[
               \Delta(T+\varepsilon\omega_A)
                         =K+\varepsilon Q.
\tag{69.28}\label{eq:verified-minkowski-perturbation}
\]
The complete character-by-character calculation is given in the proof of
\Cref{thm:ok-toric-zero-mass}; we apply it to the present notation.  More
precisely, that proof converts the book's squared-norm weight to
Guenancia's norm weight, identifies the interiors of the local Newton set
and its closure, and checks that the strict condition automatically makes
the local Laurent exponent holomorphic.  After those checks, Guenancia's
criterion \cite[Theorem~1.13]{Guenancia12} says on a maximal smooth chart
\(U_\sigma\) that the character \(\chi^u\) survives exactly when
\[
 u+\kappa_\sigma\in
 \operatorname{int}\bigl(m(K+\varepsilon Q)+\sigma^\vee\bigr),
\]
where \(\kappa_\sigma\) is the fixed all-ones vector in the chart lattice.
The proof cited above checks the Cartier-data signs, proves the global gluing
identity
\[
 K+\varepsilon Q
 =\bigcap_{\sigma\in\Sigma(n)}
   (K+\varepsilon Q+\sigma^\vee),
\]
and squeezes the strict conditions with shifts \(m^{-1}\kappa_\sigma\) between
inner and outer convex bodies.  The lattice Riemann-sum theorem then proves
that the following inner limit exists and that
\[
                         S_A(\varepsilon)
                         =\vol_n(K+\varepsilon Q),
\tag{69.29}\label{eq:verified-toric-rr}
\]
where \(\vol_n\) is normalized by the character lattice.  Thus in this toric
case the normalized section count has an actual limit.  The displayed equality
holds only after that limit, not as a literal finite-level section-count
identity, and no factor \(n!\) occurs.

Put \(d=\dim K\) and let \(A_K\) be its affine span. After translation,
decompose
\[
                  M_{\mathbf R}=V_K\oplus V_K^\perp,
       \qquad \dim V_K=d.
\]
Because \(K\) has nonempty relative interior in \(A_K\), it contains a
\(d\)-dimensional ball of radius \(a>0\) in \(V_K\). Since \(Q\) has
nonempty interior, it contains an \(n\)-dimensional ball of radius
\(b>0\). Therefore \(K+\varepsilon Q\) contains a product-type body
with \(d\) directions of fixed size and \(n-d\) directions of size
\(\varepsilon\). Hence
\[
                \vol_n(K+\varepsilon Q)
                       \geq c\varepsilon^{\,n-d}.
\tag{69.30}\label{eq:verified-volume-lower}
\]

Conversely, \(K\) is contained in a bounded \(d\)-dimensional box in
\(A_K\), and \(Q\) is contained in a bounded \(n\)-dimensional box.
Their Minkowski sum is contained in a box with \(d\) sides bounded
independently of \(\varepsilon\) and \(n-d\) sides bounded by
\(C\varepsilon\). Thus
\[
                \vol_n(K+\varepsilon Q)
                       \leq C'\varepsilon^{\,n-d}.
\tag{69.31}\label{eq:verified-volume-upper}
\]
Equations \eqref{eq:verified-cao-characterization}--
\eqref{eq:verified-volume-upper} show that the exponent in Cao's
criterion is \(n-d\). Therefore
\(\nu_{\mathrm{Cao}}(T)=d\).
\end{proof}

\begin{example}[Segment in a surface]
\label{ex:verified-toric-segment}
If \(n=2\) and \(K\) is a nondegenerate segment, then
\[
              \vol_2(K+\varepsilon Q)\asymp\varepsilon.
\]
Thus \(\nu_{\mathrm{Cao}}(T)=1\), although the actual top
nonpluripolar mass \(2!\vol_2(K)\) is zero.
\end{example}

For a concrete normalization, take
\(K=[0,a]\times\{0\}\) and \(Q=[0,1]^2\). Then
\[
             K+\varepsilon Q=[0,a+\varepsilon]\times[0,\varepsilon],
 \qquad
             \vol_2(K+\varepsilon Q)=a\varepsilon+\varepsilon^2.
\]
The leading exponent is one, so
\eqref{eq:verified-cao-characterization} gives
\(2-\nu_{\mathrm{Cao}}(T)=1\). This computes
\(\nu_{\mathrm{Cao}}(T)=1\) without confusing numerical dimension with the
vanishing two-dimensional volume of \(K\) itself.

\begin{remark}[Rationality boundary]
\label{rem:verified-toric-rationality}
The proof displayed in this section uses toric section growth and therefore
stops at integral or rational classes.  This is a boundary of the algebraic
method, not of the theorem: the analytic fixed-current argument of
\Cref{thm:ok-toric-real-products} covers arbitrary pseudoeffective real
classes and identifies both the Cao and non-pluripolar current-level
dimensions.  It does not use class-varying rational approximation.
\end{remark}

\section{No \texorpdfstring{\(\mathcal I\)}{I}-invariant body can recover actual mass}
\label{sec:verified-mass-no-go}

\begin{theorem}\label{thm:verified-mass-no-go}
There is no rule, defined for every compact K\"ahler manifold \(X\)
and every pseudoeffective class \([\theta]\), that assigns
\[
 \mathfrak B_{X,\theta}:
 \{\varphi\in\PSH(X,\theta)\}
       \longrightarrow
 \{\text{Lebesgue-measurable subsets of }\mathbf R^N\}
\tag{69.32}\label{eq:verified-body-assignment}
\]
for a fixed integer \(N\) and satisfies both:
\begin{enumerate}
\item if \(\varphi\sim_{\mathcal I}\psi\), then
      \(\mathfrak B(\varphi)\) and \(\mathfrak B(\psi)\) agree up to
      a Lebesgue-null set;
\item for every \(\varphi\),
      \[
          \vol_N(\mathfrak B(\varphi))
             =\int_X\left\langle
                    (\theta+\ddc\varphi)^n
                       \right\rangle.
      \tag{69.33}\label{eq:verified-body-mass}
      \]
\end{enumerate}
Convexity of \(\mathfrak B(\varphi)\) is not needed for the
contradiction.
\end{theorem}

\begin{proof}
Take \(X=\mathbf P^1\) with the Fubini--Study form \(\omega\),
normalized by \(\int_X\omega=1\). Let \(K\subset\mathbf P^1\) be a
polar Cantor set carrying an atom-free probability measure \(\mu\), as
in \cite[p.~31]{Car83}.
Solve
\[
                           \mu=\omega+\ddc\varphi.
\tag{69.34}\label{eq:verified-cantor-potential}
\]
Because \(\mu\) has no atoms, every Lelong number of \(\varphi\) is
zero. Skoda integrability then gives
\(\mathcal I(m\varphi)=\mathcal O_{\mathbf P^1}\) for every
\(m>0\), and hence
\[
                              \varphi\sim_{\mathcal I}0.
\tag{69.35}\label{eq:verified-cantor-i-equivalence}
\]
The measure \(\mu\) is supported on the polar set \(K\). In bidegree
one, the nonpluripolar product is the nonpluripolar part of the current
\(\mu\); it is supported in \(K\) but puts no mass on the pluripolar
set \(K\). Hence it vanishes and
\[
               \int_X\langle\omega+\ddc\varphi\rangle=0.
\tag{69.36}\label{eq:verified-cantor-zero-mass}
\]
For the smooth zero potential,
\[
                              \int_X\omega=1.
\tag{69.37}\label{eq:verified-smooth-mass}
\]
Property \(1\) says that
\(\mathfrak B(\varphi)\) and \(\mathfrak B(0)\) have the same volume.
Property \(2\), together with
\eqref{eq:verified-cantor-zero-mass}--\eqref{eq:verified-smooth-mass},
says that these volumes are respectively \(0\) and \(1\), a
contradiction.
\end{proof}

\begin{corollary}\label{cor:verified-ideals-insufficient}
No body constructed only from the collection of multiplier ideals
\(\{\mathcal I(m\varphi)\}_{m>0}\) can recover the actual
nonpluripolar mass of every potential.
\end{corollary}

\begin{proof}
Such a construction is invariant under
\(\mathcal I\)-equivalence, so
\Cref{thm:verified-mass-no-go} applies.
\end{proof}

\begin{remark}[\(P\)-sensitive refinements are not excluded]
\label{rem:verified-p-sensitive-body}
The polar-Cantor argument rules out only
\(\mathcal I\)-invariant assignments. It does not rule out a body
depending on the full \(P\)-type. Constructing such a body remains a
separate problem.
\end{remark}

\section{Three definitions of PL plurisubharmonicity}
\label{sec:verified-pl-psh-equivalence}

Let \(K\) be complete and nontrivially valued, and let \(X\) be a good strict
\(K\)-analytic space.  For a real-valued real PL function \(h\), consider:
\begin{enumerate}
\item residue-nef, or model, semipositivity;
\item classical psh positivity, tested after analytic scalar extension on
      every analytic curve; and
\item membership in the CLD-v2 psh sheaf, generated from smooth
      Hessian-positive functions by transfinitely iterated decreasing
      sequential closure.
\end{enumerate}
These conditions have genuinely different definitions.  The next theorem
identifies them on the largest domain reached by the presently available
regularization theorem.

\begin{theorem}\label{thm:verified-pl-psh-equivalence}
Let \(h\) be as above.
\begin{enumerate}
\item On \(X\),
\[
 \text{CLD-v2 psh}
 \quad\Longrightarrow\quad
 \text{classically psh}
 \quad\Longleftrightarrow\quad
 \text{model semipositive}.
\]
\item If \(X\) is boundaryless, all three conditions are equivalent:
\[
 \text{model semipositive}
 \quad\Longleftrightarrow\quad
 \text{classically psh}
 \quad\Longleftrightarrow\quad
 \text{CLD-v2 psh}.
\tag{69.38}\label{eq:verified-pl-psh-equivalence}
\]
For general \(X\), this equivalence holds on the relative interior.
\item If \(X\) is topologically separated, any of the three conditions
implies positive \(\dprimeddprime\)-current on the relative interior, and
globally when \(X\) is boundaryless.  No converse from weak current
positivity is asserted.
\end{enumerate}
The statement requires neither smoothness nor reducedness of \(X\).
\end{theorem}

\begin{proof}
Model semipositivity is equivalent to classical psh positivity by
\cite[Theorem~4.3]{GR25Classical}, while every CLD-v2 psh function is
classically psh by the comparison in \cite[\S1.2]{GR25Classical}.  These are
the arrows valid without a boundaryless hypothesis.

Assume that \(X\) is boundaryless and that \(h\) is model semipositive.  The
local regularization theorem of \cite[Introduction and Theorem~5.3]{GR25PL},
using \cite[Proposition~7.10]{GJR26}, supplies locally Fubini--Study functions
\(F_j\) converging uniformly to \(h\) on one neighborhood.  The exact local
normal form from \cite[\S7.9]{GJR26}, after replacing each analytic unit in
the source convention by its inverse, is
\[
 F=c+\sum_{\nu=1}^{N}a_\nu
       \max_{1\leq i\leq r_\nu}\log|f_{\nu i}|,
 \qquad a_\nu\geq0,
\]
where the \(f_{\nu i}\) are analytic units.  There are no arbitrary
independent real shifts inside the maxima.

For a positive integer \(p\), replace one maximum by
\[
 M_p=\frac1p\log\!\left(
             \sum_{i=1}^{r}\rme^{p\log|f_i|}\right).
\]
The real log-sum-exp function is smooth, convex, and nondecreasing in every
coordinate.  Since each \(\log|f_i|\) is CLD-smooth and pluriharmonic, the CLD
chain rule makes \(M_p\) smooth and Hessian-positive.  The unnormalized
finite-dimensional \(\ell^p\)-norm decreases with \(p\), so \(M_p\) decreases
to \(\max_i\log|f_i|\).  Finite nonnegative linear combinations preserve both
properties.  Thus every locally Fubini--Study \(F_j\) is CLD-v2 psh, although
it need not be CLD-smooth.

It remains to justify that a locally uniform limit is legal in the
decreasingly generated CLD sheaf.  After passing to a subsequence, arrange
\(\lVert F_j-h\rVert_\infty\leq e_j:=4^{-j}\), and set
\[
 c_j:=\sum_{q=j}^{\infty}(e_q+e_{q+1}).
\]
Then \(F_j+c_j\) decreases to \(h\).  Closure of the CLD-v2 sheaf under
decreasing sequences proves that \(h\) is CLD-v2 psh.  This establishes the
missing arrow and hence \eqref{eq:verified-pl-psh-equivalence}.  Applying the
same argument to \(\operatorname{Int}(X)\) proves the relative-interior
statement.

Finally, \cite[Theorem~7.14]{GJR26} says directly that a semipositive real PL
function on a separated good strict space has positive current on the
relative interior, globally in the boundaryless case; it has no
pure-dimensionality hypothesis.  The cited current theorem is one-way.  Its
converse would require pulling the singular current back to the analytic
curves in the curve criterion, an operation not presently available with the
needed positivity in this generality.
\end{proof}

\begin{example}[The PL equivalence does not mean smoothness]
On a strict annulus containing the unit Gauss point,
\[
                         h=\max\{0,\log|T|\}
\]
is locally Fubini--Study and hence satisfies all three conditions in
\Cref{thm:verified-pl-psh-equivalence}.  On the skeleton it is
\(\max\{0,s\}\), so it is not differentiable at \(s=0\) and is not CLD-smooth.
The log-sum-exp argument above is therefore essential.
\end{example}

\begin{remark}[Exact remaining arrows]
At a relative-boundary point, the available model-to-LFS regularization does
not prove the reverse CLD arrow.  For arbitrary non-PL functions, CLD-v2 psh
still implies classically psh, but the converse is not established.  The
trivially valued and non-strict variants also lie outside the cited theorem.
\end{remark}

\section{The boundary ledger}
\label{sec:verified-boundary-ledger}

For clarity, the following questions remain outside the theorem text:
\begin{enumerate}
\item the general trace Dirichlet problem needs extension in a fixed
      class with strict positivity;
\item trace on a singular unibranch ambient space needs equisingular
      regularization and multiplier ideals;
\item a general zero-mass Okounkov body needs genuinely
      \(P\)-sensitive data;
\item the unrestricted nontrivially valued Monge--Amp\`ere equation
      still carries support and characteristic restrictions;
\item global CLD--BFJ psh equivalence needs regularization and
      local-to-global positivity;
\item intrinsic non-Archimedean K\"ahler/Kodaira theory needs a viable
      differential positivity notion;
\item the global BJ Radon--Riesz property lacks strong compactness;
\item ordinary compact-K\"ahler YTD still needs the model/ordinary and
      automorphism bridges; and
\item the relation between class-level numerical dimension and the Newton
      bodies of individual singular currents remains distinct from the
      current-level identity of \Cref{thm:ok-toric-real-products}.
\end{enumerate}

\begin{remark}[How to read the ledger]
\label{rem:verified-read-ledger}
The table is not a conjecture list with presumed positive answers.
Some entries may require new definitions; others may fail. The only
claims promoted in this chapter are
\Cref{thm:verified-cld-criterion,thm:verified-fixed-dual-complex,%
thm:verified-trace-factorization,thm:verified-split-trace,%
thm:verified-toric-dimension,thm:verified-mass-no-go,%
thm:verified-pl-psh-equivalence}, together with the explicit cross-fan
counterexample.
\end{remark}

\section{Dependency graph}
\label{sec:verified-dependencies}

The proofs above are independent in the following sense:
\[
\begin{tikzcd}[column sep=large,row sep=large]
\text{current topology}
 \arrow[r]
&\text{CLD criterion}\\
\text{polyhedral slopes}
 \arrow[r]
&\text{cross-fan counterexample}\\
\text{strong compactness}
 \arrow[r]
&\text{fixed-complex theorem}\\
\text{trace comparison and uniqueness}
 \arrow[r]
&\text{\(\mathcal I\)-input/\(P\)-output trace map}\\
\text{toric convex mass and Cao products}
 \arrow[r]
&\text{numerical dimension}\\
\text{polar Cantor potential}
 \arrow[r]
&\text{mass-body no-go}\\
\text{LFS regularization and log-sum-exp}
 \arrow[r]
&\text{PL psh equivalence}.
\end{tikzcd}
\tag{69.39}\label{eq:verified-dependency-graph}
\]
No arrow passes through a problem listed in
\Cref{sec:verified-boundary-ledger}.

\section*{Source notes}

The CLD current operations and terminology use \cite{CLD25}; the
approximation criterion and cross-fan audit are the R2-validated
outcomes CLD--3 and CLD--4. The weak and \(d_\infty\) topologies on
\(X^{\lin}\), the exhaustion \(T_\omega\), and Question~11.16 use
\cite{BJ22}; only the fixed-dual-complex compactness
argument is promoted. Trace comparison and output uniqueness are from
\cite{DX24} and \Cref{chap:operations-trace}.  The product calculation was
R2-audited at the weaker \(\mathcal I\)-level; its present \(P\)-target is the
direct application of the sourced uniqueness theorem, so the obsolete audit
is not cited as a receipt for the updated formulation.  The toric
Newton-body dictionary is from the supplied complex manuscript.  The
real-class theorem in \Cref{thm:ok-toric-real-products} combines Cao's
current-level product \cite[Definitions~2--3]{Cao14}, gentle
quasi-equisingular approximation and toric \(\mathcal I\)-goodness from
\cite{XiaBook}, the local complex/real identity of \cite{CGSZ19}, and the
non-pluripolar calculus of \cite{BEGZ10}.  The second proof in this chapter
uses \cite[Proposition~4.4]{Cao14}, Guenancia's multiplier-ideal criterion
\cite[Theorem~1.13]{Guenancia12}, the toric character decomposition
\cite[\S4.3]{CLS11}, rational toric asymptotic Riemann--Roch, and elementary
convex geometry.
The polar-Cantor potential uses \cite[p.~31]{Car83}, as in the supplied
comparison chapter. No unresolved R1 attempt is
stated here as a theorem.  The new PL equivalence combines the published
arrows in \cite{GR25PL,GR25Classical,GJR26} with the log-sum-exp and
summable-shift proof above.  Its corrected Pass~R3 Rethlas audit returned
\texttt{correct} with zero critical errors and zero gaps; the receipt and
cryptographic digests are recorded in the editorial verification ledger.
The publication receipt is
\texttt{cld\_gr\_psh\_comparison\_r3.json}, with statement digest
\texttt{50c0c6b21d6a\ldots} and proof digest
\texttt{d44d6b4dc9b0\ldots}; the verifier was
\texttt{gpt-5.6-sol} at \texttt{xhigh} scrutiny on 14 August 2026. No
first-draft theorem in this chapter is enlarged beyond its audited
signature; the separately identified trace reconstruction is governed by
its source signature.
